\section{Discussão}
\color{black}

\subsection{A Regra Fundamental: Pipeline $>$ Modelo}

O resultado central deste trabalho corrobora e estende os achados de~\cite{Khamaisi2025}: em AAD com
dados escassos e restrições de Edge Computing, a \textbf{qualidade do pipeline de processamento supera
a complexidade do modelo}. O XGBoost com Super-Vetor bem construído supera o Tiny-AST --- Transformer
de última geração --- nas métricas DCASE-alinhadas, ao custo de $1/12$ da memória e $1/6$ da latência.
O erro NMF implementa o conceito do Autoencoder em $<0{,}5$~MB e $<5$~ms, eliminando a necessidade da
arquitetura pesada sem penalidade de performance.

\subsection{Por que Deep Learning foi Descartado como Motor}

Todos os modelos de Deep Learning testados foram descartados ou rebaixados a função auxiliar por ao
menos um dos critérios:
\begin{enumerate}
    \item \textbf{Latência:} 45--120~ms, violando o limite de 50~ms para Edge.
    \item \textbf{Memória:} 3--12~MB, violando o limite de 4~MB para TinyML.
    \item \textbf{Overfitting ao domínio:} Com dados escassos, redes profundas memorizam
    características do microfone (sensor-specific features) em vez das da falha mecânica, causando
    colapso em domain shift~\cite{Harada2023}.
\end{enumerate}

\subsection{Por que pAUC@0.1 Muda o Ranking}

A substituição de F1-Score e AUC global por pAUC@FPR~0.1 é uma das contribuições metodológicas mais
relevantes. Em sistemas de alarme industrial, a taxa de falso positivo deve ser minimizada --- um
sistema que gera 10~alarmes falsos por turno perde confiança operacional em semanas. O pAUC@0.1 mede
exclusivamente a performance na região crítica. Este deslocamento de métrica alterou o ranking: Tiny-AST
(excelente em AUC global, pAUC~0,91) caiu para 3º; XGBoost+HHT (ótimo na região de baixo FPR,
pAUC~0,94) ascendeu ao 1º.

\subsection{Protocolo LOSO e Anti-Leakage}

A maioria dos trabalhos de AAD na literatura usa split aleatório ou k-fold padrão, que permitem
vazamento de informação entre sensores/domínios. O protocolo LOSO proposto garante que o modelo nunca
viu dados do sensor de teste, medindo generalização real --- condição necessária para deployment
industrial seguro~\cite{Nishida2025}.

\subsection{Limite Operacional Documentado}

O Protocolo~B falha no cenário MIMII Slider com SNR~$-6$~dB (pAUC~0,79~$<$~0,80). Este é o
\textbf{limite operacional formal}: em ambientes com SNR~$<-3$~dB sem rótulos de anomalia,
o Protocolo~A é obrigatório ou deve-se coletar exemplos de falha.

\color{black}
